\section{Introducción}

Esta guía describe el procedimiento que debe seguir un técnico de Nezu para instalar, configurar, validar y entregar HomePilot. HomePilot es una consola local de control residencial que integra dispositivos, cámaras, asistentes, tableros y rutinas bajo una experiencia unificada.

\subsection{Objetivo}
El objetivo no es reemplazar automáticamente cada plataforma del cliente. HomePilot centraliza el control y conserva una arquitectura abierta: puede conectarse a un Home Assistant existente, desplegar uno complementario cuando no exista, o iniciar un entorno nativo limitado cuando el alcance del proyecto lo justifique.

\subsection{Alcance del instalador}
El instalador configura la infraestructura, valida red y servicios, crea la primera cuenta administrativa y enlaza las integraciones autorizadas. El cliente final recibe una consola funcional; no necesita administrar Docker, variables de entorno ni credenciales técnicas.

\subsection{Arquitectura de referencia}
\begin{center}
\begin{tikzpicture}[node distance=1.25cm, >=Latex, every node/.style={font=\small}]
\node[draw=nezuorange, rounded corners, fill=nezuorange!10, minimum width=3.1cm, minimum height=0.9cm] (client) {Cliente y dispositivos};
\node[draw=nezuorange, rounded corners, fill=black!3, right=of client, minimum width=3.1cm, minimum height=0.9cm] (hp) {HomePilot};
\node[draw=nezuorange, rounded corners, fill=black!3, right=of hp, minimum width=3.4cm, minimum height=0.9cm] (ha) {Home Assistant opcional};
\draw[->,thick,nezuorange] (client) -- (hp);
\draw[->,thick,nezuorange] (hp) -- node[above]{Bridge autorizado} (ha);
\end{tikzpicture}
\end{center}

\begin{tcolorbox}[title=Principio operativo]
HomePilot mantiene su propio inventario, espacios, usuarios y tableros. Cuando se usa Home Assistant, este actúa como puente de integración para descubrir y controlar los dispositivos que el cliente ya tiene configurados.
\end{tcolorbox}
